\begin{statement}{A.23}{Open}{The remaining ramified higher images in Task A.6}
The exact images and all fibers of
$\operatorname{Sec}(N/N'')$ and
$\operatorname{Sec}(N/\gamma_3(I_{\rm mark}))$ in
$\operatorname{Sec}(N/\gamma_3N)$ remain undetermined here
for general ramified classes. A.20 determines the entire
unramified family with its fibers. A.21 proves the required
cofinal comparison. A.22 computes the metabelian coefficient
lattices, their derived limits, and the complete next lifting
fiber.
\end{statement}
Statements A.17--A.18 exclude every class-two lift in the
mixed-inertia families of A.5--A.6 from both higher images,
including after finite extensions. A.19 now determines the
entire local metabelian image at every prime different from
three: nontrivial tame inertia belongs to the full actual cusp
packet, including its central coordinates, and zero abelian
tame inertia forces all tame inertia to vanish.
The remaining global metabelian problem is therefore the
simultaneous lifting of the specified local packets with the
global arithmetic relations and restricted convergence.
For the marked quotient, B.12--B.14 determine the entire
local image and its higher fibers, and A.32 gives the global
image on the density-$5/9$ prime set. The remaining global
problem retains the simultaneous arithmetic relations.
One of its explicit next-layer tests is the class-three
tame relation defect of a specified ramified class-two section.
After choosing class-three lifts of its local Frobenius and
inertia values, this defect lies in the central lattice $M_3$.
Dividing by $q$ gives its residue in $M_3(-1)$; central changes
alter it by $D_{3,q}$. The required corrections must be integral
and restricted. Computing the coefficient lattice does not
evaluate these defects. The further origin cochains and the
simultaneous ramified lift choices remain to be computed.

